% P15-U1 / N-x-008: product measures and iterated integration.
\chapter{Product Measures and Iterated Integration}
\label{sec:p15-product-measures}

A product model combines two measurable spaces. Its measure must agree
with the prescribed products on rectangles, and an iterated calculation
must agree with integration on that product space. The order of integration
can be changed for nonnegative functions, including those with infinite
integral, and for absolutely integrable signed functions. Mere existence
of the two iterated signed integrals does not suffice.

This treatment follows Section 1.7 of Durrett's author version of
\emph{Probability: Theory and Examples}, dated 11 January 2019
\cite[pp.~37--41]{DurrettPTEAuthor2019}. It supplies the product construction,
the section lemmas, both branches of his Fubini theorem, the finite-product
consequence and his four boundary examples. We call its nonnegative branch
\emph{Tonelli's theorem} and its integrable signed branch \emph{Fubini's
theorem}. The two are parts of the same source result. The final two
corollaries also solve the source's Exercises 1.7.1 and 1.7.2.
The measure-extension and integral foundations are those of
Chapters~\ref{sec:p15-measure-extension}
and~\ref{sec:p15-integration-expectation}.

\section{Product events, sections and the rectangle assignment}
\label{sec:p15-product-events}

\begin{hbdefinition}{def:p15-product-events}{The product sigma-algebra and sections}
Let \((X,\mathcal A)\) and \((Y,\mathcal B)\) be measurable spaces,
with nonempty underlying sets. Define
\[
 \mathcal S\defeq\{A\times B:A\in\mathcal A,\ B\in\mathcal B\},
 \qquad
 \mathcal A\times\mathcal B\defeq\sigma(\mathcal S).
\]
The members of \(\mathcal S\) are \emph{measurable rectangles}.
The expression \(\mathcal A\times\mathcal B\) has Durrett's typed
meaning here: the sigma-algebra on \(X\times Y\) generated by those
rectangles, rather than the set of ordered pairs of sigma-algebra members.
It is also commonly written \(\mathcal A\otimes\mathcal B\).
All product measurability in this chapter refers to this sigma-algebra.

For \(E\subseteq X\times Y\), its two kinds of \emph{section} are
\[
 E_x\defeq\{y\in Y:(x,y)\in E\},\qquad
 E^y\defeq\{x\in X:(x,y)\in E\}.
\]
For a function \(f\) on \(X\times Y\), write
\(f_x(y)=f(x,y)\) and \(f^y(x)=f(x,y)\).
The coordinate projections are \(\pi_X(x,y)=x\) and
\(\pi_Y(x,y)=y\).
\index{product sigma-algebra}\index{section!measurable set}
\end{hbdefinition}

\begin{hbtheorem}{thm:p15-product-sections}{proved=annex:p15.product-sections}{Measurable sections and maps into a product}
\par\nobreak
Every \(E\in\mathcal A\times\mathcal B\) has \(E_x\in\mathcal B\)
for every \(x\in X\) and \(E^y\in\mathcal A\) for every \(y\in Y\).
If \(f\) is product measurable and real or extended real valued,
every function \(f_x\) and \(f^y\) is measurable.

The product sigma-algebra is the smallest one making both coordinate
projections measurable. A map \(T:\Omega\to X\times Y\) from
a measurable space is measurable exactly when \(\pi_X\circ T\)
and \(\pi_Y\circ T\) are measurable. In particular, for positive
integers \(r,s\), under the coordinate identification,
\[
 \mathcal B(\mathbb R^r)\times\mathcal B(\mathbb R^s)
                       =\mathcal B(\mathbb R^{r+s}).
\]
These conclusions use no measures and no finiteness hypotheses.
\end{hbtheorem}

\begin{hbdefinition}{def:p15-product-rectangle-mass}{Rectangle masses and the chosen product construction}
Give the two measurable spaces measures \(\mu_1,\mu_2\).
A measure is \emph{sigma-finite} when its space is a countable union
of measurable sets of finite measure. This is the measure-space case
of the finiteness convention in Definition~\ref{def:p15-semialgebra-premeasure}.
On measurable rectangles prescribe
\[
 m(A\times B)\defeq\mu_1(A)\mu_2(B),
\]
using the nonnegative extended convention \(0\cdot\infty=0\).
The theorem below extends this assignment first to the algebra of finite
disjoint unions of rectangles, and then by the countable-cover
outer-measure construction of Theorem~\ref{thm:p15-premeasure-extension}
to \(\mathcal A\times\mathcal B\). Denote the measure obtained by
that specified construction by \(\mu_1\times\mu_2\), and also by
\(\rho\) within this chapter. The symbol is defined by this construction
even when uniqueness has not been asserted.
\index{product measure}\index{rectangle!measure}
\end{hbdefinition}

\begin{hbtheorem}{thm:p15-product-existence}{proved=annex:p15.product-existence}{Existence and sigma-finite uniqueness of product measure}
\par\nobreak
The rectangle assignment in Definition~\ref{def:p15-product-rectangle-mass}
is well defined and extends by the stated construction to a measure
\(\rho\) on \(\mathcal A\times\mathcal B\), satisfying
\[
 \rho(A\times B)=\mu_1(A)\mu_2(B)
       \qquad(A\in\mathcal A,\ B\in\mathcal B).
\]
This existence assertion does not require sigma-finiteness. If both
\(\mu_1\) and \(\mu_2\) are sigma-finite, then \(\rho\) is
sigma-finite and is the unique measure on this sigma-algebra with the
displayed rectangle values. In particular, the product of two probability
measures is a probability measure. The uniqueness conclusion with
sigma-finite factors is Durrett's Theorem 1.7.1; the unrestricted existence
conclusion uses the earlier, explicitly proved extension theorem.
\end{hbtheorem}

\begin{hbtheorem}{thm:p15-finite-product-measures}{proved=annex:p15.finite-product-measures}{Finite products and Euclidean volume}
For \(n\geq1\), let \((X_i,\mathcal F_i,\mu_i)\), \(1\leq i\leq n\),
be sigma-finite measure spaces. There is a unique measure on the
sigma-algebra generated by all rectangles \(\prod_{i=1}^n A_i\),
\(A_i\in\mathcal F_i\), such that
\[
 (\mu_1\times\cdots\times\mu_n)
       \left(\prod_{i=1}^n A_i\right)=\prod_{i=1}^n\mu_i(A_i).
\]
It is sigma-finite. Iterated binary products, under their coordinate
identifications, give the same sigma-algebra and measure regardless of
parenthesization. Permuting coordinates carries this measure to the
product with factors permuted in the same way.

For the constructed Borel Lebesgue measure \(\lambda\) on each copy
of \(\mathbb R\), this gives a Borel measure \(\lambda^n\) on
\(\mathbb R^n\) with
\[
 \lambda^n\left(\prod_{i=1}^n(a_i,b_i]\right)
                               =\prod_{i=1}^n(b_i-a_i)
       \qquad(a_i<b_i\text{ finite}).
\]
Here \(\lambda^n\) denotes the product measure, called Lebesgue
measure on the Borel sets of \(\mathbb R^n\).
\end{hbtheorem}

\section{Section integrals and exchange of order}
\label{sec:p15-tonelli-fubini}

For the rest of the theorems in this section, both factor measures are
sigma-finite and \(\rho=\mu_1\times\mu_2\).

\begin{hblemma}{thm:p15-section-measures}{proved=annex:p15.section-measures}{Measurability and integration of section measures}
Let both factor measures be sigma-finite and let
\(\rho=\mu_1\times\mu_2\). For every product measurable set \(E\), the functions
\(x\mapsto\mu_2(E_x)\) and \(y\mapsto\mu_1(E^y)\)
are nonnegative extended measurable functions, and
\[
 \int_X\mu_2(E_x)\,d\mu_1(x)
   =\rho(E)
   =\int_Y\mu_1(E^y)\,d\mu_2(y).
\]
In particular, \(\rho(E)=0\) exactly when \(E_x\) is
\(\mu_2\)-null for \(\mu_1\)-almost every \(x\), and exactly
when \(E^y\) is \(\mu_1\)-null for \(\mu_2\)-almost every \(y\).
\end{hblemma}

\begin{hbtheorem}{thm:p15-tonelli}{proved=annex:p15.tonelli}{Tonelli: nonnegative iterated integrals}
Let both factor measures be sigma-finite, let
\(\rho=\mu_1\times\mu_2\), and let
\(f:X\times Y\to[0,\infty]\) be product measurable.
The section integrals
\[
 T_1f(x)\defeq\int_Y f(x,y)\,d\mu_2(y),\qquad
 T_2f(y)\defeq\int_X f(x,y)\,d\mu_1(x)
\]
exist at every point in \([0,\infty]\), are measurable as functions
of the remaining variable, and satisfy
\[
 \int_X T_1f\,d\mu_1
       =\int_{X\times Y}f\,d\rho
       =\int_Y T_2f\,d\mu_2.
\]
Every value in this equality may be infinity; finiteness of the integral
is unnecessary. If two such product measurable functions agree
\(\rho\)-a.e., their first section integrals agree
\(\mu_1\)-a.e. and their second section integrals agree
\(\mu_2\)-a.e.
\index{Tonelli's theorem}
\end{hbtheorem}

\begin{hbtheorem}{thm:p15-fubini}{proved=annex:p15.fubini}{Fubini: absolutely integrable signed functions}
Let both factor measures be sigma-finite and let
\(\rho=\mu_1\times\mu_2\). Let
\(f:X\times Y\to[-\infty,\infty]\) be product measurable
and satisfy \(\int|f|\,d\rho<\infty\). Define
\[
 H_1(x)=\int_Y|f(x,y)|\,d\mu_2(y),\qquad
 H_2(y)=\int_X|f(x,y)|\,d\mu_1(x),
\]
and the measurable null sets \(N_1=\{H_1=\infty\}\subseteq X\),
\(N_2=\{H_2=\infty\}\subseteq Y\). For every \(x\notin N_1\),
the section \(f_x\) is integrable; for every \(y\notin N_2\),
the section \(f^y\) is integrable. Define real functions everywhere by
\[
 \begin{aligned}
 I_1f(x)&=\begin{cases}
  \int_Y f(x,y)\,d\mu_2(y),&x\notin N_1,\\
  0,&x\in N_1,
 \end{cases}\\
 I_2f(y)&=\begin{cases}
  \int_X f(x,y)\,d\mu_1(x),&y\notin N_2,\\
  0,&y\in N_2.
 \end{cases}
 \end{aligned}
\]
These functions are measurable and integrable, and
\[
 \int_X I_1f\,d\mu_1
       =\int_{X\times Y}f\,d\rho
       =\int_Y I_2f\,d\mu_2.
\]
Their absolute integrals are each at most \(\int|f|\,d\rho\).
Changing a product measurable \(f\) on a measurable \(\rho\)-null
set, while keeping it measurable, leaves both section-integral functions
unchanged a.e. in their respective factor spaces. Other measurable
finite assignments on \(N_1,N_2\) give the same outer integrals.
\index{Fubini's theorem}\index{iterated integral!exceptional sections}
\end{hbtheorem}

\begin{hbexample}{def:p15-fubini-exceptional-section}{An integrable product function with a nonintegrable section}
Let \(X=(0,1)\) with uniform measure and \(Y=\mathbb R\) with
Borel Lebesgue measure. Fix \(x_0\in(0,1)\), and let
\(f(x,y)=\mathbf1_{\{x_0\}}(x)\). Its support is the measurable
rectangle \(\{x_0\}\times\mathbb R\), whose product measure is
\(0\cdot\infty=0\). Thus \(\int|f|\,d\rho=0\), while
\(\int_Y|f(x_0,y)|\,d\lambda(y)=\infty\).
All other first sections vanish. The exceptional set is exactly
\(N_1=\{x_0\}\), and the function \(I_1f\) in Fubini's theorem
is zero everywhere. This calculation explains why the theorem's
integrability conclusion is almost everywhere rather than at every section.
\end{hbexample}

\begin{hbmethod}{def:p15-iterated-integral-method}{Choosing a legitimate order of integration}
Specify the two factor sigma-algebras and measures, verify their
sigma-finiteness, and establish product measurability of the integrand.
For a nonnegative integrand, Tonelli permits either order immediately.
For a signed integrand, first apply Tonelli to its absolute value. A
finite answer in either order establishes the integrability premise for
Fubini. Then calculate signed sections wherever they are integrable and
assign a measurable finite value on the remaining null set, as in the
theorem. The equality of the resulting outer integrals is the guarantee.
Two finite answers obtained by signed cancellation alone do not establish
the absolute-integrability premise.
\end{hbmethod}

\begin{hbexample}{def:p15-triangle-integral}{An iterated calculation on the unit square}
On two copies of the constructed uniform space \((0,1)\), put
\(f(x,y)=x\mathbf1_{\{y<x\}}\). The coordinates are measurable,
\(\{y<x\}\) is Borel, and \(f\geq0\), so Tonelli applies.
Integrating first with respect to \(y\) gives \(x\cdot x=x^2\),
then integrating over \(x\) gives \(1/3\). In the opposite order,
\[
 \int_y^1 x\,d\lambda(x)=\frac{1-y^2}{2},\qquad
 \int_0^1\frac{1-y^2}{2}\,d\lambda(y)
                   =\frac12-\frac16=\frac13.
\]
The one-variable integral formula follows by upper and lower sums on
interval partitions; the corresponding elementary integral rule is
justified in the annex with the failure calculations below. Thus the
two computations agree with \(\int f\,d(\mathbb U\times\mathbb U)\).
The strict boundary \(y<x\) is retained in the calculation, with
singleton endpoints having zero one-dimensional measure.
\end{hbexample}

\section{Failures of the exchange hypotheses}
\label{sec:p15-product-failures}

\begin{hbtheorem}{thm:p15-fubini-failures}{proved=annex:p15.fubini-failures}{Four source counterexamples and their exact causes}
The following examples are Durrett's Examples 1.7.5--1.7.8, with their
measurability and existence qualifications made explicit.
\begin{enumerate}
\item On \(\mathbb N_+^2\) with counting measure in each factor, set
 \(f(m,m)=1\), \(f(m+1,m)=-1\) for \(m\geq1\), and set
 \(f=0\) otherwise. Then
 \[
  \sum_{m\geq1}\sum_{n\geq1}f(m,n)=1,
  \qquad
  \sum_{n\geq1}\sum_{m\geq1}f(m,n)=0.
 \]
 Every section has finite support, but both product part integrals
 \(\int f^+\) and \(\int f^-\) are infinite. The product signed
 integral is undefined.
\item On \((0,1)\times(1,\infty)\), with Borel Lebesgue measures,
 let \(f(x,y)=e^{-xy}-2e^{-2xy}\). All individual sections are
 integrable and both iterated signed integrals are finite, but
 \[
 \begin{aligned}
 \int_0^1\left(\int_1^\infty f(x,y)\,d\lambda(y)\right)d\lambda(x)
   &=\int_0^1\frac{e^{-x}-e^{-2x}}x\,d\lambda(x)>0,\\
 \int_1^\infty\left(\int_0^1 f(x,y)\,d\lambda(x)\right)d\lambda(y)
   &=\int_1^\infty\frac{e^{-2y}-e^{-y}}y\,d\lambda(y)<0.
 \end{aligned}
 \]
 Here too, each product part integral is infinite.
\item Let \(X=(0,1)\) with its Borel sigma-algebra and uniform
 Lebesgue measure \(\mu_1\), and \(Y=(0,1)\) with all subsets measurable and
 counting measure \(\mu_2\). The diagonal indicator \(f(x,y)=\mathbf1_{\{x=y\}}\)
 is product measurable and nonnegative, yet
 \[
  \int_X\left(\int_Yf(x,y)\,d\mu_2(y)\right)d\mu_1(x)=1,
  \qquad
  \int_Y\left(\int_Xf(x,y)\,d\mu_1(x)\right)d\mu_2(y)=0.
 \]
 The counting measure on this uncountable set is not sigma-finite.
 No measure can make both of these iterated values equal to its integral
 of this same indicator.
\item Assume additionally the axiom of choice and the continuum
 hypothesis, in the form that a bijection
 \(b:\omega_1\to(0,1)\) exists, where \(\omega_1\) is the set
 of countable ordinals with its usual order. Transport that order by
 \(b\), writing \(b(\alpha)<'b(\beta)\) when \(\alpha<\beta\).
 On two uniform Borel copies of \((0,1)\), with measures
 \(\mu_1,\mu_2\), set
 \(f(x,y)=\mathbf1_{\{x<'y\}}\). Every section is Borel measurable,
 but
 \[
  \int_X f(x,y)\,d\mu_1(x)=0\quad\hbox{for every }y,
  \qquad
  \int_Y f(x,y)\,d\mu_2(y)=1\quad\hbox{for every }x.
 \]
 Consequently \(f\) is not product measurable. This last example is
 conditional on the stated set-theoretic assumptions; it is not an
 unconditional construction of such an order.
\end{enumerate}
\end{hbtheorem}

\section{Absolute summation and area under a graph}
\label{sec:p15-product-corollaries}

\begin{hbcorollary}{thm:p15-absolute-iterated-series}{proved=annex:p15.absolute-iterated-series}{An iterated integrability criterion and exchange of a series}
For sigma-finite factors and a product measurable extended real \(f\),
\[
 \int_X\left(\int_Y|f(x,y)|\,d\mu_2(y)\right)d\mu_1(x)<\infty
\]
is equivalent to \(\int|f|\,d\rho<\infty\), and also to the
same finite bound with the order reversed. Any of them permits the
signed equalities of Fubini's theorem, with its a.e. section convention.

Let \((Y,\mathcal B,\nu)\) be sigma-finite and let real measurable
\(f_n:Y\to\mathbb R\), \(n\geq1\), satisfy
\(\sum_n\int_Y|f_n|\,d\nu<\infty\). Then \(\sum_n f_n\)
converges absolutely a.e. Define \(S\) to be that sum where it
converges absolutely and zero on the measurable null exceptional set.
The function \(S\) is integrable and
\[
 \int_Y S\,d\nu=\sum_{n\geq1}\int_Y f_n\,d\nu,
 \qquad
 \sum_{n\geq1}\left|\int_Y f_n\,d\nu\right|<\infty.
\]
\end{hbcorollary}

\begin{hbcorollary}{thm:p15-layer-cake}{proved=annex:p15.layer-cake}{Area under a measurable graph and the tail integral}
Let \((X,\mathcal A,\mu)\) be sigma-finite and let
\(g:X\to[0,\infty]\) be measurable. In
\(X\times[0,\infty)\), give the second factor its Borel
Lebesgue measure and define
\(D_g=\{(x,t):0\leq t<g(x)\}\). This set is product measurable and
\[
 \int_X g\,d\mu
   =(\mu\times\lambda)(D_g)
   =\int_{[0,\infty)}\mu(\{x:g(x)>t\})\,d\lambda(t).
\]
The value may be infinite. In particular, for a nonnegative random
variable \(Z\) on a probability space,
\(\mathbb E[Z]=\int_0^\infty\mathbb P(Z>t)\,d\lambda(t)\).
\index{layer-cake formula}\index{expectation!tail integral}
\end{hbcorollary}

\begin{hbwarning}{def:p15-product-completion-domain}{The sigma-algebra remains part of the theorem}
Every section-measurability statement above concerns the product
sigma-algebra generated by measurable rectangles. No completion of
either factor or of the product measure has been silently substituted.
Passing to a completed product domain requires its own measurable
representative and exceptional-section argument. In particular, the
fact that a function has measurable sections does not prove its joint
measurability, as the conditional fourth counterexample demonstrates.
\end{hbwarning}

\begin{proofboundary}{def:p15-product-next-boundary}{Further product and probability constructions}
These results construct finite products and justify the stated iterated
integrals and countable-series interchange. Durrett's remaining
Section 1.7 exercises on Stieltjes integration by parts, translated
cumulative functions and the sine integral are separate applications.
Independence is developed in
\cref{sec:p15-independent-random-variables}; densities, convolutions
and the gamma and normal sum laws are developed in
\cref{sec:p15-convolution-gamma-normal}.
\Cref{sec:p15-countable-product-measures} constructs countable products
from consistent finite-dimensional laws and extends that construction
through explicit Borel real codings.
\Cref{sec:p15-radon-nikodym,sec:p15-conditional-expectation,sec:p15-regular-conditional-laws}
supply density representation, conditional expectation and regular kernels.
A general completion theorem for product measures requires its further
representative and exceptional-section construction.
\end{proofboundary}
